\section{Dataset Construction}
\label{sec:dataset-construction}

To evaluate the translation and to validate the output empirically, we constructed a dataset drawn from real world open-source Elixir projects.
The dataset records extracted TypeSpec annotations, translated Elixir Type annotations, function definition, associated user-defined (or remote) types, and outcomes of Dialyzer and the new typechecker from the compiler.
This section describes how the dataset is built, how each tool is invoked, and how the results are structured.

\subsection{Overview}
% And the repos (projects)?
% Remove type_examples later.
The dataset covers \textbf{383 repositories from 54 distinct owners} (GitHub organisations and individual maintainers, with three repositories hosted on other forges) and contains \textbf{\num{23073} translated function specifications}.
Construction proceeds in three sequential stages driven by three Mix tasks:
\begin{enumerate}
    \item \textit{Translation execution} (\texttt{mix run\_translation}), which translates every \texttt{@spec} in the corpus and emits \texttt{dataset.jsonl}.
    \item \textit{Dialyzer analysis} (\texttt{mix run\_dialyzer}), which runs Dialyzer on the translated projects and enriches each dataset entry with Dialyzer outcomes.
    \item \textit{Typechecker execution} (\texttt{mix run\_typecheck}), which runs the new Elixir compiler's typechecker and enriches entries with its outcomes.
\end{enumerate}
Each stage reads and rewrites the same target dataset file so that the three stages compose into a single artefact.

\subsection{Dataset Entry Schema}

Each line of \texttt{dataset.jsonl} is a JSON object with the following fields:

\begin{itemize}
    \item \texttt{project}, \texttt{module}, \texttt{function}, \texttt{arity}: the identity of the function.
    \item \texttt{file}: the source file path relative to the corpus root (e.g.\ \texttt{cldr/cldr47/lib/cldr.ex}). Because the same \texttt{(module, function, arity)} may occur in several sub-projects under one organisation, this field is what disambiguates a warning's true target; see Subsection~\ref{sec:warning-association}.
    \item \texttt{spec}: the original \texttt{@spec} annotation as a string.
    \item \texttt{elixir\_type}: the translated annotation produced by the migrator.
    \item \texttt{type}: the list of \texttt{@type} definitions in scope for the function.
    \item \texttt{definition}: the full function source reconstructed from the AST.
    \item \texttt{spec\_lines}, \texttt{start\_line}, \texttt{end\_line}: source location metadata.
    \item \texttt{dialyzer}: an object with fields \texttt{analysis} (\texttt{"success"} or \texttt{"failed"}), \texttt{pass} (Boolean), and \texttt{messages}. Each message records its \texttt{category}, \texttt{line}, \texttt{message} text, and an explicit \texttt{target} (the \texttt{Module.function/arity} the warning is attributed to).
    \item \texttt{typecheck}: an object with fields \texttt{analysis}, \texttt{pass}, and \texttt{messages}, each message recording \texttt{target}, \texttt{line}, and \texttt{message}.
\end{itemize}

\subsection{Stage 1: Translation Execution}

Translation is orchestrated by \texttt{TranslationRunner.main/1}.
It begins with a full-corpus type-cache build: all source files are passed to \texttt{Migrator.TypeTranslator.process/1}, which scans every \texttt{@type}, \texttt{@opaque}, and \texttt{@typep} declaration across the repository and builds a module-indexed table of parsed type definitions.
This global pre-pass is necessary so that user-defined types are available for substitution when TypeSpec annotations from individual files are translated later.

Files are then processed concurrently using \texttt{Task.async\_stream} with one task per source file and a concurrency limit of \texttt{System.schedulers\_online()}, matching the number of logical CPU cores available at runtime.
Each task is independent: it reads the source file, extracts specs, invokes the migrator, and produces a list of dataset entries without communicating with other tasks.
Each task proceeds as follows.

\paragraph{TypeSpec extraction.}
The file is parsed into an AST and walked by a recursive extractor that traverses \texttt{defmodule}, \texttt{\_\_block\_\_}, \texttt{@spec}, \texttt{def}, and \texttt{defp} nodes.
For each function that has a matching \texttt{@spec}, the extractor records a structured entry containing the spec text, input and output AST nodes, guard clauses, source line numbers, and the full function body reconstructed via \texttt{Macro.to\_string/1}.
End-line positions are derived by walking the function's AST with \texttt{Macro.prewalk} and tracking the maximum line number seen.

\paragraph{Elixir Type injection.}
The migrator is invoked with the \texttt{:descr\_assert} output mode on the pre-modification file, producing a list of translated annotations keyed by the \texttt{@spec} line numbers.
This is done before any file writes so that the line numbers used by the extractor and by the migrator are consistent.
The migrator also inserts \texttt{@assert\_type\_form} lines on a line above corresponding TypeSpecs in the target source file; these are later consumed by the typechecker.

\paragraph{Struct representation and its resolution limit.}
\label{sec:struct-representation}
A struct type is translated to a set-theoretic map keyed on \texttt{\_\_struct\_\_}.
Because every Elixir module is represented on the BEAM by an \texttt{Elixir.}-prefixed atom---a full struct name is an atom, for example, a struct defined in a module named \texttt{Some} is fully expanded to \texttt{:"Elixir.Module"}~\cite{elixir-modules-atoms}; a bare \texttt{:"Plausible.Site"} is a different atom that a \texttt{\%Plausible.Site\{\}} pattern can never match.
When the struct's fields are all resolvable, the type is emitted as a \emph{closed} map, faithful to the fact that a struct admits no unregistered keys.
When the defining module lies outside the translation path, so its \texttt{defstruct} is unresolved, the exact closed type cannot be written, hence, the struct is switched to a open map and then intersected with (\texttt{dynamic()}) retaining \texttt{\_\_struct\_\_}-keyed field, consistent with the translator's treatment of every other unresolvable definition (a remote type outside the path is likewise replaced with \texttt{dynamic()}).
Although switching to an open map is broadening, this is necessary because a struct type likely has its fields and since there is no way to find out the filed types, we choose a type that is compatible rather than leaving it incorrect.
% Two residuals remain and are irreducible without a symbol table: a struct referenced through a local alias the translator cannot expand to its full module path, and a struct whose field set cannot be recovered.
% These are the genuine cost of AST-based translation without a persistent lookup table (PLT), and their downstream effect on the typechecker is quantified in Subsection~\ref{sec:typecheck-only}.

\paragraph{Annotation matching.}
Translated annotations are joined to spec entries by matching on the \texttt{@spec} line number.
Each entry also records its \texttt{file} (the source path relative to the corpus root), which later disambiguates warnings whose \texttt{(module, function, arity)} is not unique across the organisation's sub-projects.
The result for each file is a list of entries ready for serialisation.

\paragraph{Fault tolerance.}
Each file task is wrapped in \texttt{with\_retry/2}, which re-attempts the task up to three times on exception and logs a warning on each failure.
Tasks that exceed a 60-second timeout are killed and omitted from the output.

The entries from all tasks are collected and written to \texttt{results/dataset.jsonl} as newline-delimited JSON.

\paragraph{DatasetWriter.}
For runs that accumulate results incrementally, a \texttt{DatasetWriter} GenServer provides concurrent write access to the file.
Write requests arrive as asynchronous \texttt{cast} messages so that callers are never blocked.
Before each write, the server loads the existing file, computes a deduplication key of \texttt{(project, module, function, arity)} for each incoming entry, removes any existing entries with matching keys, and then writes the merged content atomically.
A synchronous \texttt{flush} call, used at shutdown, drains the message queue before the process terminates.

\subsection{Stage 2: Dialyzer Analysis}

Dialyzer analysis is performed by \texttt{DialyzerRunner.main/1}.
The runner loads the existing dataset to determine which projects are present, then discovers all Mix project roots under the translation output directory by globbing for \texttt{mix.exs} files and filtering out umbrella sub-applications to avoid double-counting.

\paragraph{Dependency injection.}
Before running Dialyzer, the runner checks whether \texttt{dialyxir} is already listed as a dependency in the project's \texttt{mix.exs}.
If it is absent, the runner injects \texttt{\{:dialyxir, "\textasciitilde> 1.4", only: [:dev, :test], runtime: false\}} into the dependency list using a regular-expression replacement.
This allows Dialyzer to be invoked through \texttt{mix dialyzer} regardless of the original project configuration.

\paragraph{Project processing.}
Each project is processed concurrently.
The runner invokes \texttt{mix deps.get}, \texttt{mix compile}, and \texttt{mix dialyzer} in sequence using \texttt{System.cmd/3}.
If dependency resolution or Dialyzer itself fails, the project is marked as skipped and its entries receive \texttt{"analysis":"failed"} in the dataset.

\paragraph{Warning extraction.}
Dialyzer's output is parsed in two passes to handle the two output formats produced by different versions of dialyxir.
The first pass detects legacy-format blocks introduced by the string \texttt{"Legacy warning:"} and normalises them.
The second pass operates on the structured section beginning at \texttt{"Total errors:"}.
Within each section, warning blocks are delimited by lines matching the pattern \texttt{lib/*.ex:N:} and are parsed with a regular expression that extracts the file path, line number, warning category, and message text.
The modern dialyxir format tags each warning with a structured category (\texttt{no\_return}, \texttt{call}, \texttt{pattern\_match}, and so on); the legacy format does not, so for legacy blocks the category is inferred from the message text and the full message is preserved (Subsection~\ref{sec:warning-association} details this normalisation).

Because Dialyzer reports only a file and line, never the enclosing function, each warning is mapped back to a \texttt{(module, function, arity)} target by re-parsing the source file and locating the spec'd clause whose line range contains the warning line.
The warning's file is recorded relative to the corpus root so that it can be matched against the entry's \texttt{file} field.

\paragraph{Dataset enrichment.}
After all projects have been processed, the runner iterates through the base entries and, for each entry, finds the Dialyzer warnings whose \texttt{(file, module, function, arity)} match the entry's identity.
The \texttt{file} component is essential: an organisation such as \texttt{cldr} ships several release versions as separate sub-projects, each defining the same \texttt{Cldr.\,*} functions, and a key omitting the file would attach a warning from one version to all of them.
An entry passes if no matching warnings exist.
The enriched entries are written back to \texttt{dataset.jsonl}.

\subsection{Stage 3: Typechecker Execution}

Typechecker execution is performed by \texttt{TypecheckRunner.main/2}.
The key difference from the Dialyzer stage is that the new Elixir compiler's typechecker is invoked via a custom Elixir binary built from the development branch of the next Elixir version (1.21.0-dev) rather than through a standard Mix plugin, since the current stable 1.19.* releases do not expose the typechecker.
The runner prepends the directory of the custom binary to the shell \texttt{PATH} before invoking \texttt{mix compile}, so that the typechecking compiler is used in place of the system Elixir installation.

\paragraph{Warning parsing.}
The typechecker emits warnings in the format \texttt{warning: <message>} followed by a location line beginning with a box-drawing corner prefix that records the file path, line number, and a fully qualified \texttt{Module.function/arity} identifier.
Unlike Dialyzer's line-only output, this location line names the target function authoritatively, so no source re-parsing is needed to attribute a typecheck warning.
Three properties of the compiler's output must be handled for the parse to be complete and correct, and are detailed in Subsection~\ref{sec:warning-association}: most warning headers are \emph{indented}, the location line carries an optional column and an occasional \texttt{(app version)} prefix, and the human-readable message is stored separately from the authoritative target.

\paragraph{Arity mismatch.}
Although Dialyzer allows omitting an optional argument's type, the new compiler does not support such behaviour yet, hence, it mandates including an optional argument's type and fails when a spcification does not contain the full argument types, evoking an arity mismatch error.
To increase quantity of our dataset, we visited every airty failure and added a \texttt{dynamic()} type as a filler.

\paragraph{Dataset enrichment.}
Matched warnings are joined to dataset entries by the \texttt{(file, module, function, arity)} key, following the same file-disambiguated pattern as the Dialyzer stage.
The final \texttt{dataset.jsonl} contains complete records with both tool outcomes, each carrying its explicit \texttt{target}, for every translated function.

\subsection{Warning Association}
\label{sec:warning-association}

The scientific value of the dataset rests on every warning being attributed to the function that actually produced it, since the downstream training signal pairs a function definition with the verdicts on that definition.
Two properties of the corpus make naïve attribution unsafe, and the construction addresses both.

\paragraph{File-disambiguated keys.}
The \texttt{project} field denotes a GitHub organisation, and within one organisation the same \texttt{(module, function, arity)} frequently recurs: the \texttt{cldr} organisation alone ships four release versions as independent Mix projects, each defining \texttt{Cldr.validate\_script/1} and many others.
Across the corpus, 195 \texttt{(project, module, function, arity)} keys are shared by two or more distinct source files (\num{683} entries in total).
Keying warning attribution on identity alone would broadcast a warning emitted by one version to every same-named entry.
Both stages therefore key on \texttt{(file, module, function, arity)}, where each warning's file is normalised to the same corpus-root-relative form as the entries'.
The relative path is invariant between the translated and non-translated source trees, so a Dialyzer warning (computed on the non-translated tree) and a typecheck warning (computed on the translated tree) both match the single entry whose \texttt{file} they share.

\paragraph{Target provenance.}
The two tools locate the target differently, and the dataset records the result explicitly in each message's \texttt{target} field.
Dialyzer emits only \texttt{file:line}; its target is \emph{inferred} as the spec'd function whose source line range contains the warning line.
The typechecker emits a \texttt{Module.function/arity} location line; its target is taken authoritatively from that line.
In both cases the message text often names a \emph{callee} (e.g.\ a function the body calls incorrectly) rather than the enclosing function; recording an explicit \texttt{target} alongside the message text removes any ambiguity for downstream consumers about which definition a verdict describes.

%% This is a convenient tool that is used for cleaning up annotations in some cases.
% \subsection{Annotation Cleanup}

% Before the analysis stages run, the \texttt{AnnotationCleaner} module removes all \texttt{@assert\_type\_form} lines that were injected into source files during translation.
% It applies a line-filter to every \texttt{.ex} file under the project directory, excluding the migrator's own source tree, ensuring that the analysis tools operate on clean source without the injected annotations interfering with compilation or warning attribution.

%% This also helps to obserb the gathered data, but not essentially neccessray for data construction.
% \subsection{Dataset Analysis Report}

% The \texttt{DatasetParser} module reads the completed \texttt{dataset.jsonl} and generates a self-contained \LaTeX\ report.
% Entries are classified into four groups: those where both tools passed, those where both warned, those where only Dialyzer warned, and those where only the new typechecker warned.

% For entries that triggered warnings, a rule-based classifier assigns a human-readable category and explanation.
% Dialyzer warnings are classified into categories such as opaqueness violation, no-local-return, call-will-not-succeed, contract violation, pattern-match coverage, and unknown function, based on the warning category tag and the presence of keywords in the message text.
% Typechecker warnings are classified as return-type mismatch, incompatible argument types, or general type inconsistency.
% Entries containing \texttt{dynamic()} in the translated annotation receive a separate sub-category indicating that annotation imprecision may have masked or caused the warning.

% The report includes a summary table with per-project breakdowns of how many specs were successfully analysed and how many triggered warnings from each tool, a warning-category index grouped by frequency, and per-entry listings showing the original spec, translated annotation, function body, and tool messages side-by-side.
